% Options for packages loaded elsewhere
\PassOptionsToPackage{unicode}{hyperref}
\PassOptionsToPackage{hyphens}{url}
\PassOptionsToPackage{dvipsnames,svgnames,x11names}{xcolor}
\documentclass[
  11pt,
]{article}
\usepackage{xcolor}
\usepackage[margin=0.85in]{geometry}
\usepackage{amsmath,amssymb}
\setcounter{secnumdepth}{-\maxdimen} % remove section numbering
\usepackage{iftex}
\ifPDFTeX
  \usepackage[T1]{fontenc}
  \usepackage[utf8]{inputenc}
  \usepackage{textcomp} % provide euro and other symbols
\else % if luatex or xetex
  \usepackage{unicode-math} % this also loads fontspec
  \defaultfontfeatures{Scale=MatchLowercase}
  \defaultfontfeatures[\rmfamily]{Ligatures=TeX,Scale=1}
\fi
\usepackage{lmodern}
\ifPDFTeX\else
  % xetex/luatex font selection
\fi
% Use upquote if available, for straight quotes in verbatim environments
\IfFileExists{upquote.sty}{\usepackage{upquote}}{}
\IfFileExists{microtype.sty}{% use microtype if available
  \usepackage[]{microtype}
  \UseMicrotypeSet[protrusion]{basicmath} % disable protrusion for tt fonts
}{}
\makeatletter
\@ifundefined{KOMAClassName}{% if non-KOMA class
  \IfFileExists{parskip.sty}{%
    \usepackage{parskip}
  }{% else
    \setlength{\parindent}{0pt}
    \setlength{\parskip}{6pt plus 2pt minus 1pt}}
}{% if KOMA class
  \KOMAoptions{parskip=half}}
\makeatother
\usepackage{longtable,booktabs,array}
\newcounter{none} % for unnumbered tables
\usepackage{calc} % for calculating minipage widths
% Correct order of tables after \paragraph or \subparagraph
\usepackage{etoolbox}
\makeatletter
\patchcmd\longtable{\par}{\if@noskipsec\mbox{}\fi\par}{}{}
\makeatother
% Allow footnotes in longtable head/foot
\IfFileExists{footnotehyper.sty}{\usepackage{footnotehyper}}{\usepackage{footnote}}
\makesavenoteenv{longtable}
\setlength{\emergencystretch}{3em} % prevent overfull lines
\providecommand{\tightlist}{%
  \setlength{\itemsep}{0pt}\setlength{\parskip}{0pt}}
\usepackage{bookmark}
\IfFileExists{xurl.sty}{\usepackage{xurl}}{} % add URL line breaks if available
\urlstyle{same}
% fallback for those not using the hyperref driver hyperxmp:
\makeatletter
\@ifundefined{xmpquote}{\newcommand{\xmpquote}[1]{#1}}{}
\makeatother
\hypersetup{
  pdftitle={Refining Section 4.1 of the original draft},
  colorlinks=true,
  linkcolor={black},
  filecolor={Maroon},
  citecolor={Blue},
  urlcolor={blue},
  pdfcreator={LaTeX via pandoc}}

\title{Refining Section 4.1 of the original draft}
\usepackage{etoolbox}
\makeatletter
\providecommand{\subtitle}[1]{% add subtitle to \maketitle
  \apptocmd{\@title}{\par {\large #1 \par}}{}{}
}
\makeatother
\subtitle{Stochastic regional masking: replacement text and verified
explanation}
\author{}
\date{8 October 2026}

\makeatletter
\renewcommand{\maketitle}{\noindent{\LARGE\@title\par}\vspace{0.6em}
\noindent{\small\@date\par}\vspace{1em}}
\makeatother

\begin{document}
\maketitle

This note supplies a replacement for Section 4.1 of
\texttt{wacv2007-v1.pdf} and explains the evidence behind it. The
top-level notebooks remain the source of truth. Implementation details
below come from the functions those notebooks call; new diagnostic
calculations are identified separately. The original PDF and notebooks
have been preserved.

\subsection{Replacement text for Section
4.1}\label{replacement-text-for-section-4.1}

\subsubsection{4.1. Stochastic regional
masking}\label{stochastic-regional-masking}

During training, S-JEPA predicts numerical descriptions of selected
parts of a skeleton sequence using the remaining observations. A
\emph{mask} identifies the positions to withhold from the context
network. Each 32-frame window includes all 33 body landmarks. Four
successive frames for one landmark form a token, giving eight time
blocks and 264 tokens per window. The twelve landmarks described below
guide how masks are sampled; the model continues to receive the full set
of landmark positions before masking.

A fresh mask is drawn for each training example from seven overlapping
body regions: head, mouth, left arm, right arm, trunk, left leg and
right leg. Each draw chooses a region, a span of one to six consecutive
time blocks, and a start position that keeps the span inside the window.
Span length and the possible start positions are sampled uniformly. The
selected positions are added to the mask until at least 158 tokens have
been selected, or 36 region draws have been attempted. Because the last
region is added in full, the final hidden fraction can exceed the
nominal target of 60\%.

The sampling weights reflect an anatomical design choice motivated by
gait analysis. Regions containing either shoulder, hip, knee, ankle,
heel or foot index receive a multiplier of 1.5. The weight is then
divided by the average number of regions containing the region's
landmarks, before the weights are converted to selection probabilities.
This adjustment partly accounts for shoulders and hips belonging to two
regions. Arm regions receive the multiplier through their shoulders, so
elbows, wrists and hand landmarks also inherit that region weight. These
rules do not give every highlighted landmark the same hiding frequency,
or make its final frequency 1.5 times that of other landmarks.

The sampler retains at least one visible token from the twelve-landmark
set somewhere in the window. That token may belong to a shoulder, so the
rule does not ensure visible hips or legs in every example. Overlapping
draws can hide a joint for all eight time blocks, and a whole time block
can be hidden when context remains elsewhere. In notebook 02's separate
bank of 512 masks generated with seed zero, the mean hidden fraction was
63.3\%. Every landmark was visible in at least one time block in at
least 74.0\% of those sampled windows. Availability as context varied
across time: for example, left-hip context occurred in 51.6\% of samples
at the first block but 7.0\% at the fifth block.

These observations describe the sampling rule, rather than the masks
encountered throughout model training. Clinical gait studies motivate
attention to body position and coordination, but do not validate this
particular twelve-landmark choice or its weighting factor. The study
includes no controlled comparison with uniform, motion-dependent or
alternative anatomical masking. Its results therefore leave the
contribution of this masking design to classification performance
unresolved.

\clearpage
\subsection{What required correction, and
why}\label{what-required-correction-and-why}

The original section already recognized that the twelve landmarks were
not a fixed mask and that the masking design had not been evaluated
separately. Those qualifications should be retained. Its main weaknesses
were an incomplete sampling rule, an ambiguous summary of coverage, and
insufficient separation between clinical motivation and evidence of
benefit. The visible editorial marker ``Rewrite:'' should also be
removed.

\textbf{The multiplier alone does not describe the sampler.} The
implementation applies both the 1.5 anatomical multiplier and a
correction for overlapping region membership. Omitting the second
operation prevents an accurate reconstruction of the procedure. It also
obscures why trunk draws have a lower probability than head draws, even
though the trunk contains selected landmarks. The probabilities below
apply to one region draw; they are not final joint-hiding probabilities.

{\def\LTcaptype{none} % do not increment counter
\begin{longtable}[]{@{}llr@{}}
\toprule\noalign{}
Region & Landmark indices & Probability per draw \\
\midrule\noalign{}
\endhead
\bottomrule\noalign{}
\endlastfoot
Head & 0--8 & 12.7854\% \\
Mouth & 9, 10 & 12.7854\% \\
Left arm & 11, 13, 15, 17, 19, 21 & 16.4384\% \\
Right arm & 12, 14, 16, 18, 20, 22 & 16.4384\% \\
Trunk & 11, 12, 23, 24 & 9.5890\% \\
Left leg & 23, 25, 27, 29, 31 & 15.9817\% \\
Right leg & 24, 26, 28, 30, 32 & 15.9817\% \\
\end{longtable}
}

The exact rule is \(w_G=b_G/m_G\), followed by \(P(G)=w_G/\sum_H w_H\).
Here \(b_G\) is 1.5 when region \(G\) contains a selected landmark and 1
otherwise; \(m_G\) is the mean number of region memberships among its
landmarks. The twelve selected indices are 11, 12, and 23--32. The
factor 1.5 is a configured design choice; the retained experiment does
not establish that it is optimal.

\textbf{``Approximately 60\%'' conceals useful implementation details.}
The code rounds \(0.6\times264\) to 158 and adds complete region spans
without trimming the last draw to that budget. This explains the
diagnostic average of about 63.3\%. The one-to-six-block limit applies
separately to each region draw. Several overlapping draws can cover all
eight blocks for the same joint.

\textbf{Coverage has two different denominators.} ``Visible somewhere in
a window'' counts a window once if a joint has any context token.
``Visible token fraction'' counts individual joint--time positions. The
former can look reassuring while a joint is rarely available at
particular times. Notebook 02 reports both measures, including the low
middle-block hip coverage. Its minimum 74.0\% window coverage is an
observation from 512 sampled masks, not a guarantee for each window,
each batch or a saved training run. The guarantee in the code concerns
one selected-set token somewhere in each window. In the first seed-zero
demonstration, both complete legs are hidden for the entire window.

\textbf{The original 5,000-mask summary should not be presented as
notebook output.} Notebook 02 saves a 512-mask diagnostic. Its printed
mean is rounded to 0.63; the replacement's 63.3\% comes from
recalculating that same seeded bank, which gives 0.6327607. A separately
reproduced 5,000-mask, seed-zero calculation yields 77.84--78.31\%
hiding for shoulders, 76.87--77.26\% for hips, 62.06--62.54\% for the
other leg landmarks, and 63.61--64.30\% for elbows and wrists. For
facial landmarks, head-region tokens are hidden 59.65\% of the time and
mouth tokens 52.18\%. Their joint-weighted average is 58.29\%. Thus the
original ``58\% for facial landmarks'' is compatible with an aggregate
across all eleven facial landmarks, but it conceals different rates for
the two regions. The replacement uses the notebook's own bank and keeps
the expanded calculation in this verification note.

\textbf{Clinical relevance supplies a rationale, not a measured
benefit.} Notebook 02 maps landmarks to candidate gait measurements, but
the table does not show that those physical measurements were recovered
or clinically validated in this experiment. Its introductory suggestion
that motion-based masking would conceal the relevant disease signal is
also an untested explanation: withheld information is used as a
prediction target, so withholding alone does not demonstrate loss of
useful learning. The paper can state that this implementation uses
anatomical weights without claiming that motion-dependent masking is
unsuitable for MS or PD.

\subsection{Related wording that should remain
consistent}\label{related-wording-that-should-remain-consistent}

If the original landmark illustration is retained, replace its caption
with:

\emph{The highlighted landmarks are both shoulders, hips, knees, ankles,
heels and foot indices (BlazePose indices 11, 12 and 23--32). They
influence the weights used to sample body regions during training. All
33 landmarks remain part of the input, and a new regional mask is drawn
for each example. Weights include both a factor of 1.5 for regions
containing highlighted landmarks and a correction for overlapping region
membership. Highlighting therefore indicates the sampling design, rather
than a fixed hidden set or a measured contribution to classification
accuracy.}

The reflection limitation belongs beside the augmentation description or
in the Discussion, with a cross-reference if needed. Reflection swaps
left and right landmark slots in the context input, while the target
mask and unreflected reference retain their original indices.
Consequently, a physical joint selected as a reference target can appear
in the opposite-side context slot. This is a mismatch within training
examples; it is not evidence of overlap with held-out test recordings.
The retained experiments do not measure its effect on accuracy.
Correcting the text does not produce results for a corrected training
procedure.

The classification readout also needs separate wording: the fixed
seed-zero mask selects 161 \textbf{output} positions from the target
encoder after it receives the complete window. It does not hide those
input coordinates during classification. Confusing that readout with the
changing training masks would give readers the wrong account of how
predictions are made.

\clearpage
\subsection{Evidence and verification}\label{evidence-and-verification}

Notebook cell indices here are zero-based and include Markdown cells.
Links are relative to this note's folder.

{\def\LTcaptype{none} % do not increment counter
\begin{longtable}[]{@{}
  >{\raggedright\arraybackslash}p{(\linewidth - 2\tabcolsep) * \real{0.5000}}
  >{\raggedright\arraybackslash}p{(\linewidth - 2\tabcolsep) * \real{0.5000}}@{}}
\toprule\noalign{}
\begin{minipage}[b]{\linewidth}\raggedright
Evidence
\end{minipage} & \begin{minipage}[b]{\linewidth}\raggedright
What it establishes
\end{minipage} \\
\midrule\noalign{}
\endhead
\bottomrule\noalign{}
\endlastfoot
\href{../../02_anatomical_mask_and_tokenization.ipynb}{Notebook 02},
cells 7, 9, 10--13 and 16--17 & Token count; twelve indices; 512-mask
saved diagnostic; coverage definitions and first-mask demonstration. \\
\href{../../sjepa/masking_v2.py}{Mask implementation},
\texttt{JOINT\_GROUPS}, \texttt{\_region\_weights},
\texttt{sample\_target\_mask} & Seven regions, membership correction,
draw probabilities, budget, span lengths, stopping rule and context
guarantee. \\
\href{evidence/notebook-audit.md}{Notebook audit} and
\href{evidence/notebook-audit-checks.json}{stored checks} & Independent
deterministic 5,000-mask diagnostic; separate classification readout and
reflection correspondence checks. \\
\href{evidence/section-4.1-verification.md}{Section-specific
verification} & Rechecked cell/line mapping and numerical reproduction
for this standalone refinement. \\
\href{evidence/venue-and-literature.md}{Literature audit} &
Primary-source support and limits for retained scientific claims. \\
\end{longtable}
}

Abdelfattah and Alahi's
\href{https://www.ecva.net/papers/eccv_2024/papers_ECCV/papers/04755.pdf}{S-JEPA
paper (ECCV 2024)} supplies the relevant skeleton-prediction reference.
Its motion-aware rule increases the sampling probability of joints with
larger motion; it should not be described as always selecting the
highest-motion joints. That comparison is background to this
implementation, rather than an experiment performed here.
\href{https://www.nature.com/articles/s41598-018-22676-0}{Filli et
al.~(2018)} and
\href{https://www.nature.com/articles/s41598-020-68263-0}{Plotnik et
al.~(2020)} support gait-related clinical motivation without validating
the selected mask. The rewritten section therefore needs no claim that
the present design is clinically established or superior to those
alternatives.

The numerical mask checks require no pose re-extraction, model
retraining or change to the notebooks. They verify what the documented
sampler does. Determining whether the anatomical weights improve
learning would require a new, matched experiment.

\end{document}
